% ═══════════════════════════════════════════════════════════════
% SPANISCH LEHRERHINWEISE TEMPLATE
% ═══════════════════════════════════════════════════════════════
% Fachfarbe: #c41e3a (Karminrot)
%
% Enthält:
% - Kurzplan (Stundenverlauf)
% - Differenzierungshinweise [B]/[S]/[E]
% - Häufige Fehler + Reaktionen
% - Material-Checkliste
% - Lösungen
% ═══════════════════════════════════════════════════════════════

\documentclass[11pt, a4paper]{article}

% ═══════════════════════════════════════════════════════════════
% SW-MODUS TOGGLE
% ═══════════════════════════════════════════════════════════════
% Setze \swmodustrue VOR \begin{document} für Schwarz-Weiß-Druck
\newif\ifswmodus
\swmodusfalse    % Standard: Farbe (auf true setzen für SW)

% ═══════════════════════════════════════════════════════════════
% PAKETE
% ═══════════════════════════════════════════════════════════════

\usepackage[utf8]{inputenc}
\usepackage[T1]{fontenc}
\usepackage[ngerman]{babel}
\usepackage{helvet}
\renewcommand{\familydefault}{\sfdefault}

\usepackage[margin=2cm]{geometry}
\usepackage{enumitem}
\usepackage{tabularx}
\usepackage{booktabs}
\usepackage{longtable}

\usepackage{xcolor}
\usepackage{amssymb}            % Für \square (Checkboxen)
\usepackage{tcolorbox}
\tcbuselibrary{skins}

\usepackage{fancyhdr}
\usepackage{lastpage}
\usepackage{needspace}          % Für \needspace (Seitenumbruch-Kontrolle)

% ═══════════════════════════════════════════════════════════════
% FARBEN
% ═══════════════════════════════════════════════════════════════

% Basis-Farben
\definecolor{spanisch-color}{HTML}{c41e3a}
\definecolor{grau-color}{HTML}{888888}
\definecolor{hellgrau-color}{HTML}{f5f5f5}
\definecolor{basis-color}{HTML}{2a7a4b}
\definecolor{standard-color}{HTML}{2c5aa0}
\definecolor{erweiterung-color}{HTML}{7b1fa2}
\definecolor{loesung-color}{HTML}{c62828}

% SW-Farben (druckerfreundlich: keine Füllflächen, nur Linien)
\definecolor{sw-dark}{HTML}{000000}
\definecolor{sw-gray}{HTML}{666666}
\definecolor{sw-white}{HTML}{ffffff}

% Bedingte Farbzuweisung
\ifswmodus
    \colorlet{spanisch}{sw-dark}       % Rahmen/Text: schwarz
    \colorlet{grau}{sw-gray}           % Sekundärtext: dunkelgrau
    \colorlet{hellgrau}{sw-white}      % Hintergrund: weiß (keine Füllung)
    \colorlet{basis}{sw-dark}          % [B] Marker: schwarz
    \colorlet{standard}{sw-dark}       % [S] Marker: schwarz
    \colorlet{erweiterung}{sw-dark}    % [E] Marker: schwarz
    \colorlet{loesung}{sw-dark}        % Lösungen: schwarz
\else
    \colorlet{spanisch}{spanisch-color}
    \colorlet{grau}{grau-color}
    \colorlet{hellgrau}{hellgrau-color}
    \colorlet{basis}{basis-color}
    \colorlet{standard}{standard-color}
    \colorlet{erweiterung}{erweiterung-color}
    \colorlet{loesung}{loesung-color}
\fi

\newcommand{\fachfarbe}{spanisch}

% ═══════════════════════════════════════════════════════════════
% VARIABLEN
% ═══════════════════════════════════════════════════════════════

\newcommand{\thema}{[THEMA]}
\newcommand{\sequenz}{[SEQ]}
\newcommand{\block}{[BLOCK]}
\newcommand{\fach}{Spanisch}
\newcommand{\klasse}{7}
\newcommand{\dauer}{90 Min. (Doppelstunde)}
\newcommand{\lerngruppengroesse}{24}
\newcommand{\datum}{\today}

% ═══════════════════════════════════════════════════════════════
% KOPF- UND FUSSZEILE
% ═══════════════════════════════════════════════════════════════

\pagestyle{fancy}
\fancyhf{}
\renewcommand{\headrulewidth}{0.4pt}
\renewcommand{\footrulewidth}{0.4pt}

\lhead{\fach{} -- Klasse \klasse}
\chead{\textbf{LH-\sequenz\block{} (Lehrerhinweise)}}
\rhead{\datum}

\lfoot{}
\cfoot{}
\rfoot{Seite \thepage\ von \pageref{LastPage}}

% ═══════════════════════════════════════════════════════════════
% BEFEHLE
% ═══════════════════════════════════════════════════════════════

% Differenzierungsmarker
\newcommand{\B}{\textcolor{basis}{\textbf{[B]}}}
\newcommand{\Std}{\textcolor{standard}{\textbf{[S]}}}
\newcommand{\E}{\textcolor{erweiterung}{\textbf{[E]}}}

% Lösungen
\newcommand{\loesung}[1]{\textcolor{loesung}{\textbf{#1}}}

% Spanisch-Hervorhebung
\newcommand{\es}[1]{\textcolor{\fachfarbe}{\textbf{#1}}}

% Boxen (SW-kompatibel: schwebender Titel auf weißem Grund)
\newtcolorbox{kurzplan}{
    colback=hellgrau,
    colframe=\fachfarbe,
    colbacktitle=white,
    coltitle=\fachfarbe,
    boxrule=1pt,
    arc=0pt,
    fonttitle=\bfseries,
    attach boxed title to top left={yshift=-2mm, xshift=5mm},
    boxed title style={boxrule=0pt, colback=white},
    title=Kurzplan
}

\newtcolorbox{differenzierung}{
    colback=white,
    colframe=grau,
    colbacktitle=white,
    coltitle=grau,
    boxrule=0.5pt,
    arc=0pt,
    fonttitle=\bfseries,
    attach boxed title to top left={yshift=-2mm, xshift=5mm},
    boxed title style={boxrule=0pt, colback=white},
    title=Differenzierung
}

\newtcolorbox{fehlerbox}{
    colback=white,
    colframe=loesung,
    colbacktitle=white,
    coltitle=loesung,
    boxrule=0.5pt,
    arc=0pt,
    fonttitle=\bfseries,
    attach boxed title to top left={yshift=-2mm, xshift=5mm},
    boxed title style={boxrule=0pt, colback=white},
    title=Häufige Fehler
}

\newtcolorbox{materialbox}{
    colback=white,
    colframe=\fachfarbe,
    colbacktitle=white,
    coltitle=\fachfarbe,
    boxrule=0.5pt,
    arc=0pt,
    fonttitle=\bfseries,
    attach boxed title to top left={yshift=-2mm, xshift=5mm},
    boxed title style={boxrule=0pt, colback=white},
    title=Material-Checkliste
}

% ═══════════════════════════════════════════════════════════════
% DOKUMENT
% ═══════════════════════════════════════════════════════════════

\begin{document}

% ═══════════════════════════════════════════════════════════════
% KOPFBEREICH
% ═══════════════════════════════════════════════════════════════

\begin{center}
    {\Large\bfseries\textcolor{\fachfarbe}{Lehrerhinweise: \thema}}\\[0.3em]
    {\normalsize LH-\sequenz\block{} | \dauer}
\end{center}

\vspace{10pt}

% ═══════════════════════════════════════════════════════════════
% 1. KURZPLAN
% ═══════════════════════════════════════════════════════════════

\section*{1. Stundenverlauf (Kurzplan)}

\textbf{1. Stunde (45 Min.)}

\begin{tabularx}{\textwidth}{|c|l|X|l|}
\hline
\textbf{Zeit} & \textbf{Phase} & \textbf{Inhalt} & \textbf{Material} \\
\hline
0--5 & Einstieg & [Beschreibung] & Tafel \\
\hline
5--15 & Input & [Beschreibung] & Audio \\
\hline
15--25 & Übung 1 & [Beschreibung] & -- \\
\hline
25--35 & Übung 2 & [Beschreibung] & AB \\
\hline
35--45 & Sicherung & [Beschreibung] & Tafel \\
\hline
\end{tabularx}

\vspace{1em}

\textbf{2. Stunde (45 Min.)}

\begin{tabularx}{\textwidth}{|c|l|X|l|}
\hline
\textbf{Zeit} & \textbf{Phase} & \textbf{Inhalt} & \textbf{Material} \\
\hline
0--5 & Wiederholung & [Beschreibung] & -- \\
\hline
5--20 & Anwendung & [Beschreibung] & Karten \\
\hline
20--35 & Festigung & [Beschreibung] & AB \\
\hline
35--45 & Abschluss & [Beschreibung], HA & Tafel \\
\hline
\end{tabularx}

% ═══════════════════════════════════════════════════════════════
% 2. DIFFERENZIERUNG
% ═══════════════════════════════════════════════════════════════

\section*{2. Differenzierung}

\begin{differenzierung}
\textbf{Legende:} \B{} Basis (BR) | \Std{} Standard (MR) | \E{} Erweiterung (GY)

\vspace{0.5em}

\textbf{WICHTIG:} Diese Marker sind NUR für die Lehrkraft. Auf Schülermaterial erscheinen KEINE Niveaumarkierungen!
\end{differenzierung}

\vspace{1em}

\begin{tabularx}{\textwidth}{|l|X|X|X|}
\hline
& \B{} \textbf{Basis} & \Std{} \textbf{Standard} & \E{} \textbf{Erweiterung} \\
\hline
\textbf{Aufgabe 1} & Zuordnung mit Hilfe & Zuordnung selbstständig & + Sätze bilden \\
\hline
\textbf{Aufgabe 2} & Lückentext mit Wortliste & Lückentext & + Eigene Beispiele \\
\hline
\textbf{Aufgabe 3} & Dialog mit Vorlage & Dialog mit Stichworten & Freier Dialog \\
\hline
\textbf{Aufgabe 4} & Optional / verkürzt & Vollständig & + Erweiterung \\
\hline
\end{tabularx}

\vspace{1em}

\textbf{Konkrete Hinweise:}
\begin{itemize}
    \item \B{} SuS mit LRS: Zeitzugabe (+10 Min.), vergrößerte AB-Version
    \item \B{} DaZ-SuS: Wortschatz deutsch-spanisch vorab besprechen
    \item \E{} Leistungsstarke: Expertenrolle bei Partnerarbeit, Zusatzaufgaben
\end{itemize}

% ═══════════════════════════════════════════════════════════════
% 3. HÄUFIGE FEHLER
% ═══════════════════════════════════════════════════════════════

\section*{3. Häufige Fehler}

\begin{fehlerbox}
\begin{tabularx}{\textwidth}{|l|X|X|}
\hline
\textbf{Fehler} & \textbf{Beispiel} & \textbf{Reaktion} \\
\hline
Aussprache [X] & *[falsch] statt \es{[richtig]} & Chorisches Sprechen, Einzelkorrektur \\
\hline
Verwechslung [A/B] & *[falsch] statt \es{[richtig]} & Kontrastive Übung, Visualisierung \\
\hline
Interferenz Deutsch & *[deutsche Struktur] & Bewusstmachung, L1-Vergleich \\
\hline
Akzente vergessen & *[ohne Akzent] & Hinweis auf Bedeutungsunterschied \\
\hline
\end{tabularx}
\end{fehlerbox}

% ═══════════════════════════════════════════════════════════════
% 4. MATERIAL-CHECKLISTE
% ═══════════════════════════════════════════════════════════════

\section*{4. Material-Checkliste}

\begin{materialbox}
\begin{itemize}[label=$\square$]
    \item AB-\sequenz\block.pdf (\lerngruppengroesse{} Kopien)
    \item ML-\sequenz\block.pdf (1× für Lehrkraft)
    \item Lehrwerk: Qué pasa 1, S. [X--Y]
    \item Audio-Track [X] (USB-Stick / Streaming)
    \item Tafelmarker / Kreide
    \item Optional: LP-\sequenz\block.html (Tablets/PCs)
    \item Optional: Bildkarten / Flashcards
\end{itemize}
\end{materialbox}

% ═══════════════════════════════════════════════════════════════
% 5. LÖSUNGEN
% ═══════════════════════════════════════════════════════════════

\section*{5. Lösungen AB-\sequenz\block}

\textbf{Aufgabe 1: Zuordnung} (4 Punkte)

\begin{tabular}{ll}
\es{[ES\_1]} & $\rightarrow$ \loesung{[DE\_X]} \\
\es{[ES\_2]} & $\rightarrow$ \loesung{[DE\_Y]} \\
\es{[ES\_3]} & $\rightarrow$ \loesung{[DE\_Z]} \\
\es{[ES\_4]} & $\rightarrow$ \loesung{[DE\_W]} \\
\end{tabular}

\vspace{1em}

\textbf{Aufgabe 2: Lückentext} (4 Punkte)

\begin{tabular}{ll}
a) & \loesung{[LÖSUNG\_A]} \\
b) & \loesung{[LÖSUNG\_B]} \\
c) & \loesung{[LÖSUNG\_C]} \\
d) & \loesung{[LÖSUNG\_D]} \\
\end{tabular}

\vspace{1em}

\textbf{Aufgabe 3: Dialog} (6 Punkte)

Mögliche Antworten:
\begin{itemize}
    \item B1: \loesung{Me llamo [Name]. / Soy [Name].}
    \item B2: \loesung{Soy de Alemania. / Soy de [Ort].}
    \item B3: \loesung{¡Igualmente! / ¡Mucho gusto también!}
\end{itemize}

\textit{Bewertung: Je 2P pro korrekte, sinnvolle Antwort}

\vspace{1em}

\textbf{Aufgabe 4: Freie Produktion} (6 Punkte)

Bewertungskriterien:
\begin{itemize}
    \item Begrüßung vorhanden: 1P
    \item Frage vorhanden: 1P
    \item Antwort passend: 2P
    \item Verabschiedung vorhanden: 1P
    \item Sprachliche Korrektheit: 1P
\end{itemize}

% ═══════════════════════════════════════════════════════════════
% 6. HAUSAUFGABE
% ═══════════════════════════════════════════════════════════════

\section*{6. Hausaufgabe}

\begin{itemize}
    \item Lehrwerk S. [X], Aufgabe [Y]
    \item Vokabeln lernen: [Liste von-bis]
    \item Optional: Lernpfad LP-\sequenz\block.html (digital)
\end{itemize}

% ═══════════════════════════════════════════════════════════════
% 7. REFLEXIONSNOTIZEN
% ═══════════════════════════════════════════════════════════════

\section*{7. Reflexionsnotizen (nach Durchführung)}

\begin{tabularx}{\textwidth}{|l|X|}
\hline
\textbf{Was lief gut?} & \\[2em]
\hline
\textbf{Was lief nicht?} & \\[2em]
\hline
\textbf{Zeitmanagement} & \\[2em]
\hline
\textbf{Für nächstes Mal} & \\[2em]
\hline
\end{tabularx}

\end{document}
